\chapter{李新野短文精選}

這裏精選李新野過去文章四篇。均刊載於李新野的個人中文網誌。分別爲《論中學生結婚》、《論張桂梅》、《我的父親李松堅》、《我的姐姐李新瑩》。尤其是前三篇，在網絡上受到熱烈歡迎。

可以訪問李新野的個人中文網誌 https://sinyalee.com/blog 和英文網誌 https://sinyalee.com/essays 閱讀更多的文章。

\newpage

\section{論中學生結婚}

一對十幾歲的少男少女，在鄉村生活之中互生情愫，心生愛戀。被雙方父母祝福，明媒正娶。按照中華傳統風俗互相拜堂結爲夫妻。浪漫，莊重之至，正常人都會心生羨慕，給予祝福。這對年輕的夫妻現在一定承受着巨大的心理壓力，希望他們能接收到我的祝福。

反觀現在的所謂都市現代男女，一羣30歲上下的大叔阿姨，找對象和菜市場買菜、豬農選種一樣，先看看對方的身高體重長相，再看看對方收入、有沒有房、有沒有車，工資卡上不上交再決定結不結婚。

再看看廣大鄉村，在墮女嬰和適齡女性大規模進城的雙重壓力之下，幾千萬光棍，爲了傳宗接代已接近瘋狂。農村剩下的女性出租子宮。男性三個錢包幾十萬彩禮，讓女性給他們生一個孩子，之後女性便離婚改嫁，孩子也不養了，繼續嫁給下一個農村光棍，給另外一家人傳宗接代。

就這麼一羣人，批判這對情侶不是「自由戀愛」，照照鏡子，看看你們配嗎？

\ 

說什麼「移風易俗」的，也要看什麼背景。在50年前，反對早婚，有其進步意義。現在已經2020了，中國已經發生了翻天覆地的變化，再來談什麼移風易俗，就讓人笑掉大牙了。

2020年，自由戀愛，婚前性行爲早就遍地都是。高中生戀愛、性行爲也十分常見，打開電視都是婚前意外懷孕人工流產的廣告。這些全部都是完全合法的行爲。而這對夫婦和其他人的唯一區別，就是被父母祝福，確定了夫妻關係。請問讓他們不結婚確定關係而是進行沒有保障的婚前性行爲、非婚生子、甚至墮胎、對身體造成不可恢復的傷害，對他們就是好事，對中華民族就是好事？

說教育的，也是完全不知所謂。要知道，廣東是發達地區，文風鼎盛。這對夫婦根本就不是什麼社會底層。中國現在高中入學率實際上只有55\%，而丈夫是高中生。也就是說這對夫婦的教育程度至少已經超過45\%的中國人！教育程度高，自由戀愛，被父母祝福，明媒正娶，爲中華民族養育有知識有文化的下一代，有什麼不好？讓某些教育程度低，窮困潦倒，大字不認識幾個的人生一堆暴恐分子就是爲中國好？

\ 

再說說我的經歷吧。我在汕頭的時候住的小區是全汕頭最好的小區之一。就我認識的鄰居裏面，有上市公司宜華木業的老闆、和榮誠食品的老闆。樓下一直停着兩輛東瀛戰神日產GTR跑車。

而我們樓裏，就有一對15歲結婚生子的夫妻。也是在初中自由戀愛，然後父母見面之後便明媒正娶。生了一個可愛的男寶寶，十分美好。沒有惡臭的相親，也沒有惡臭的巨額彩禮。他們結婚之後也繼續接受教育。因爲早婚早育，爺爺奶奶外公外婆都很年輕，有時間和精力幫忙帶孫子。而且他們也請了兩三個保姆，一點問題都沒有。估計他們結婚以後在大學裏面也能專心學習，而不是整天想着男男女女勾心鬥角失戀傷心。

我經常在電梯裏面，看到年輕的媽媽抱着兒子到樓下花園玩耍。媽媽穿着一個漂亮的白色帶花的夏裙，寶寶眼睛大大的，臉圓圓的，超級可愛。人世間最美好的畫面非此莫屬了吧。

說這件事，並不是說只有請得起兩三個保姆的有錢人才能早婚早育。我只是想要說明，不要以「移風易俗」的名義，發泄你們心中對農村人民的歧視和仇恨。希望你們這些把早婚和「社會底層」聯繫在一起的城市底層打工人能清醒一下，對廣大農村同胞多點包容，少點仇恨。

\ 

--------------------------------------------------------------------------

\ 

\textit{2020年11月寫於美國紐約，原載於知乎。知乎討論主題是當時汕頭潮陽貴嶼鎮一個18歲和一個14歲的中學生結婚的新聞。當時在互聯網上受到了很多女權分子的咒罵。這是我在知乎火力全開攻擊女權邪教的第一篇文章。每一句都是戳中國女權的肺管子（想嫁有錢人嫁不出去的濫交打胎的城市底層打工女性），導致這些邪教徒對我恨之入骨。2021-2022年是中國女權的巔峯，搞出了吳亦凡、董志民兩大冤案。而中國男性大規模覺醒，還要等到2023年底和2024年。這篇文章絕對是中國反女權邪教運動的重要的一篇文章。這篇文章一經發表，10分鐘點讚已過百。之後一個小時內被刪除。但是一時「洛陽紙貴」，很多人複製抄送。}

\textit{李新野網誌原文連結：https://sinyalee.com/blog/?p=935}

\newpage

\section{論張桂梅}

張桂梅就是一個沒生過小孩、思想極端的女拳頭子罷了。

她自己沒生過小孩，不知道養育子女的偉大和辛苦，自然不懂家庭主婦的價值。自己又斷子絕孫了，並且沒有勇氣承認自己是錯誤的，所以只能夠欺騙自己說不生小孩是好事，然後再拼命給學生宣揚不生小孩的思想，來給自己心理安慰。不然你讓她怎麼活？她採訪說什麼教育改變三代人，真可笑，真的被她洗腦的女學生都不知道有沒有三代人了。

當然了，張桂梅本身思想再極端，危害也僅限於這個學校。而宣揚張桂梅，錯誤嚴重一萬倍。不說她的絕育思想了，你們難道沒有意識到對張桂梅和「走出大山的女孩」的宣傳充滿了歧視嗎？當然，知乎用戶大部分是城市打工人，對鄉下的同胞缺乏同理心。我換個說法你就懂了：

\textit{十五年前，李新野校長放棄了美國紐約的高薪工作，回到饒平山區，創辦了饒平女子留美預備學校。十五年來，讓多少女學生走出了大山，走出了中國。李新野的學生都是來自饒平和周邊大埔、平和、詔安等山區縣。很多客家山區女孩子世世代代都只能給中國人生小孩，沒法獲得教育，沒法離開中國。而現在，李新野90\%以上的學生都拿到了美國全獎獎學金，從美國大學畢業，給猶太人老闆打工，和印度婆羅門同事談笑風生，有的還嫁給了美國白人，獲得了更好的生活。}

你覺得不舒服嗎？不舒服就對了。農村人看到對張桂梅的吹捧，也是這種感覺。一個山區外面的人，自上而下的要讓山區人離開山區，這不是地域歧視、城鄉歧視是什麼？山區女性就應該離開山區甚至嫁給城市人，而不能和山區男性一樣守護、建設自己的家鄉，這不是歧視女性是什麼？而山區男性卻一定要留在家鄉，不能上張桂梅的學校，不能進城，這不是歧視男性是什麼？按照美國最先進女權主義的說法，對張桂梅的吹捧，充滿了「老白男從上而下的男性說教（mansplaining）」和「種族歧視微歧視（racial microaggression）」，還充滿了男性對女性的「性別刻板印象（gender stereotypes）」。

總而言之，張桂梅不僅僅是「過譽」了。她的錯誤的極端的思想，根本就不應該得到任何的宣傳。希望大家能正確批判張桂梅，宣傳正能量，讓城市和農村人民都能愛國愛鄉，建設祖國，建設家鄉，實現城鄉共同富裕，而不是整天想着嫌棄自己的父親、兄弟、兒時青梅竹馬的男友，整天就想着去縣城、去北上廣深、去美國。一個社會有少數極端利己主義者這麼想沒多大問題，而把這種思想當正確思想來宣傳，那就大錯特錯了。

\ 

--------------------------------------------------------------------------

\ 

\textit{2020年12月寫於美國紐約，言論自由受美國憲法第一修正案保護。原載於知乎，得到過萬點讚後被刪。}

\textit{李新野網誌原文連結：https://sinyalee.com/blog/?p=928}

\newpage

\section{我的父親李松堅}
\label{3faces}

我的父親李松堅，是一個地產商。他有一個頭、一個脖子、卻有着三張臉。

\ 

第一張臉，是對主人屈膝諂媚的臉。

和大多數人一樣，我有不超過一個父親。可是我父親不一樣。從我出生的時候開始，我父親的父親就一直改變。從他高考落榜進社會不久，就開始認爹了。只要誰能夠濫用職權，把人民羣衆的利益輸送給他，他就認誰做乾爹。（後來年齡實在太大，沒法認爹了，就改認大哥。）可能是因爲跟我父親走得太近，我父親最崇敬的父兄，很多都進去了，比如上海前書記陳良宇，湖北前書記蔣超良。

我的父親是一個表裏如一的人，他是發自內心的崇拜他的這些主人們。他從髮型，公司發言，到男女關係，都事無巨細地向他們學習。連他自己在澄海南洋鄉下的會議室，都嚴格按照上海市委的樣式裝修：長方形的大廳，兩個單人沙發在大廳的一端，背面是漂亮的背景牆。兩個沙發的中間是一個茶幾。兩側是兩排單人沙發，襯托着中間兩個位置。

他也很希望我能繼承他的遺志，一直是這麼教育我的。我電腦競賽全國第二保送清華的時候，他非常開心，送給我一本書《有官在身》。他告訴我，你情商太低，要把這本書仔仔細細地讀一遍，學習怎麼讓上級開心。我高中的時候，升學第一志願是麻省理工，第二志願是香港科大，第三志願是中山大學。他知道之後，非常生氣，嚴肅地批判我。他跟我說：「你現在這麼有才華，應該去北京上學，也許哪個首長的女兒就看上你了呢！」於是，我作爲USACO（美國電腦奧賽）全球第一名，從麻省理工轉到了清華大學。

可惜我不爭氣，在清華四年，最後還是沒有一個首長的女兒看上我。

\  

我父親的第二張臉，是對神奉獻一切的臉。

我的父親一生都在追求愛情。或許是因爲他沒文化，或許是因爲他認了太多爹，亦或許是因爲他長了一張曹德旺的臉，除了他媽（這個是親媽），從來沒有一個女人愛他。

他窮盡一生追求女性的認可。在他眼中，女人就是神。在他的公司上海明園，女性員工的福利比北歐還好：只要說自己來月經了，就可以比男性員工多放幾天的假。我父親把這個福利叫做「月假」。

我的父親最愛的女人，叫做凌菲菲。他們結婚的時候，凌菲菲是一個在廣東打工的二婚帶男孩的女人。我父親作爲一個潮汕男人，衝破了封建的藩籬，忤逆了自己的父母（這個也是親的），毅然決然地跟她走入婚姻的殿堂。

2005年，我父親跟着陳良宇一起，挪用社保基金，被調查、入獄、無法說話的時候，凌菲菲趁機毅然決然地跟他離婚，分走了他一半，也就是幾十億的資產。也許正是凌菲菲這種凌厲的作風，讓我父親感受到了女神的天罰般的疼痛爽感。在他刑滿釋放的時候，回到上海，和凌菲菲又手牽手重歸於好。那幾十億的奉獻，他也甘之若飴。只可惜凌菲菲錢到手了，不肯復婚。

以前，我覺得幾十億還是挺過分的。因爲算下來，我父親和他的女神，一次要一千萬左右。但是後來發生了一件事，讓我改變了想法。2021年的時候我回國，想要幫凌菲菲把公司做大做強，於是我去上海香港廣場的蘋果店採購蘋果電腦工作用。當時銷售說可以開一個商業賬戶，結果我報了上海明園集團的公司名和稅號，銷售說之前和明園的銷售，都由凌菲菲女士負責。我當時就釋然了——我想到2016年父親節的時候，在美國買了一個蘋果手機寄給他當禮物。當時打電話，我報快遞單號，結果他連26個英文字母都不知道怎麼寫下來——一個用蘋果電腦，收藏現代藝術的優雅女性，竟然被不懂拼音的我父親糟蹋了那麼多年，女神這種痛苦，是多少錢都補償不了的吧。他們的關係，在她心中，和在他心中，是那麼一致。

\ 

李松堅的第三張臉，是對奴才們蔑視的臉。

在他的心中，不是主子，不是神，就是奴才了。包括我在內，中國的年輕男性們，在他眼中都是他的奴才。對了，我姐也是奴才。雖然她是女的，但是不能被操，所以也不屬於女神。我姐上了大學，就去上海跟我父親在一起。她畢業後，她本來跟凌菲菲一樣，在我父親相鄰的辦公室工作。但是過了一段時間，就被我父親攆走了。

我父親很討厭高學歷的人。他很喜歡教育我說，不要以爲我考上清華就多麼厲害。「你看那些復旦交大的高材生，不也是給我打工。」

哦對，有一個例外，就是那些在學校裏有編制的教授和老師。比如我的大姨李麗麗，又比如上海大學的主任教授施利毅。上海師範大學畢業的施利毅教授的上海大學納米技術研究中心，跟我父親說他們有着先進的科學技術，於是我父親大手一揮，幾億元投資血本無歸——也許他和施主任坐在會議室正中的沙發上談話，就很快樂吧。

我父親雖然表面是一個商人，卻對商業嗤之以鼻。21年我回國，我用經濟和金融知識跟他說，房地產已經到頂了。結果他聽完，比我高中的時候說要去麻省理工的時候還要生氣。他說我在美國太久，被洗腦洗壞了，不懂中國的國情。在他的世界觀裏，西方的經濟學不適合中國，中國房子永遠漲。他的財富不是來自於商業和金融，而是神的賞賜——他作爲一個對主人、對神奉獻的虔誠信徒，得到主人和神的青睞，所以才能獲得永遠漲的上海房產，住在上海徐匯的頂層複式公寓，睥睨着幾十層腳下石庫門裏面的奴才。

「你們這些復旦交大的奴才們，就應該背三十年房貸，給我打工，讓我的房子永遠漲。」2021年，我的父親這麼想。

「誒，你們這羣奴才們，怎麼不買我的房了呢？你們怎麼這麼不努力！怎麼這麼不像話！」2025年，我的父親望着明園集團到處的爛尾樓，內心發出了怒吼。

——哦，對了，忘了說。他的公司現在一個復旦交大的人都沒有。他眼裏的奴才們，都離職了。

\ \\李新野\\
2025年6月15日\ 父親節\ 寫於曼哈頓


\ 

--------------------------------------------------------------------------

\ 

\textit{本文初版的截圖在微信瘋狂轉發，傳閱無數，一時洛陽紙貴。}

\textit{李新野網誌原文連結：https://sinyalee.com/blog/?p=1001}

\newpage

\section{我的姐姐李新瑩}
\label{sister}

在社會上行走久了，感覺很多人不像人，更像是畜生：有人是雞、有人是鴨、有人是牛馬。有的畜生向上修煉成精，比如我的父親，地產商李松堅，我媽媽很小的時候就告訴我他其實是豬精\footnote{我媽媽對李松堅的稱呼是潮汕話「豬哥精」。}。而有的畜生則向下繼續墮落，比如我姐姐李新瑩，給雞當狗，成爲畜生的畜生。

\ 

我姐姐的畜生道，還要從1992年的那個夏夜開始說起。

1989年，我的姐姐出生在廣東汕頭。本來她是拿着宗馥莉般的大女主劇本的：家族長女，爺爺是鎮長，外公是校長，奶奶和外婆是泰國和臺灣歸國家庭的千金，媽媽是體制內中學教師，而我的父親，在90年代初就是資產過億的大老闆了。可是，沒過幾年，父親就出軌了。他出軌的小三，是一個在廣東打工的二婚帶男娃的狐狸精，名字叫做凌菲菲。

很快，父親就和狐狸精同居了。爲了討狐狸精開心，他把我姐姐接過去和狐狸精同住，跟狐狸精一起睡，還讓我姐姐叫狐狸精媽媽。狐狸精有肺結核。很快，我的姐姐也被傳染了肺結核。當時我姐姐才3歲左右，身體抵抗力很差。很快她就肺結核發病，高燒，昏迷。當時醫生把我姐姐抱起來，她的頭、手都是直接耷拉下去，一點意識都沒有。醫生說，這個小女孩可能救不活了。我的父親看到我姐姐這個樣子，突然良心發現。他抱着我的姐姐，用盡所有方法搶救。用了超過成人劑量幾倍的抗生素和激素，終於把我姐姐搶救回來了。

在那個1992年的夏夜，三歲的姐姐從瀕死狀態蘇醒過來，迷迷蒙蒙之中，睜開雙眼，是父親偉岸的身影，和充滿父愛的呼喚。她又閉上了雙眼，安靜地睡去。但是那一夜那炙熱的父愛，已經永遠地植入了她心中的最深處。

但是我的姐姐不知道的是，父親對她的父愛，僅限於她差點被狐狸精害死的那一夜。很快，父親就嫌她煩人，把她送回給了我媽媽，和狐狸精雙宿雙棲去了。送回去的時候，父親並沒有把我姐姐得肺結核瀕死的事情告知我媽媽，我和我媽媽和我外公都感染了結核菌。我和媽媽抵抗力比較強，沒有發展成肺結核。但是我外公肺結核發病，留下了嚴重的後遺症，最後也因爲肺部問題去世了。至於我姐姐，在被送回來之後，身體一直很差。康復後半邊臉都是黑點，輕微毀容。而且不知道是因爲高燒昏厥、抗生素還是激素，我姐姐受到了永久性的智力損傷，成爲家族最愚笨，成績最差的一個。\footnote{30年來，李松堅和凌菲菲一直隱瞞我姐姐肺炎瀕死的事，而且也隱瞞故意把肺結核傳染給我的事情，並逃逸定居上海至今。我的姐姐和您一樣，也是從我這裏第一次知道這件事。}

\ 

豬精和狐狸精雙宿雙飛之後，轉移了所有資產，並僱傭黑社會威脅欺負我們：監聽電話、扔死老鼠、辱罵、威脅。我媽媽最後害怕黑社會，也怕我姐姐撫養權被搶走被搞死，被迫簽下離婚協議，只拿到千分之一左右的資產。

可是，偶蹄目的豬，怎麼敵得過食肉目的狐狸呢？豬精因爲陳良宇案挪用社保基金，被紀委調查的時候。狐狸精毅然決然地背刺，在豬精消失的時候，背刺起訴，把豬精靠黑社會從我媽媽身上搶奪來的資產，分走一半。包括上海明園集團的一半股份和上海的多套別墅房產。而他最愛的繼子，也不認他了。

豬精被背刺之後，一面假惺惺把我和姐姐接到上海，一面深入反思：以後再也不能找食肉動物了！豬腦子轉了一圈又一圈，還是叫雞好——雞，不僅不是食肉動物，連哺乳動物都不是！於是，他到上海音樂學院，認識了一個龜公，叫做徐沛東。徐沛東當時是上海音樂學院的教授。他在他的藝術生裏面，找了一個願意做雞的，叫做許潔。龜公哼哧哼哧把雞背在身上，獻給了豬精。爲了豬精開心，龜公還利用他教授和校董的身份和豬精的錢，賄賂學校，給許潔買了一個教授的身份，變成了教授雞。後來，教授雞在豬精接近60歲高齡的時候，給豬精生了個雄性雞仔，叫做李新權。

有了弟弟之後，我和姐姐的命運就被安排好了。因爲我不僅美貌還聰明，獲得了中國電腦奧賽第二名和美國電腦奧賽第一名，我爸就逼我放棄麻省理工的錄取轉去清華，希望我能和領導的子女和親。至於被狐狸精燒壞腦袋的姐姐，出嫁也收益不大，於是我父親把她留在身邊當下人。

2021年我回國，豬精把我和姐姐叫到公司喫飯。他爲了給自己小時候拋妻棄子和謀殺子女未遂辯護，對我們說：「你們年輕人，不應該談戀愛，不應該結婚，應該奮鬥事業。等到你們50歲左右才結婚最好。」當時我震驚得不行：50歲我和姐姐都絕經了，還結個屁婚。我驚恐地望向姐姐，結果姐姐專心致志，聽得津津有味。當時我姐姐做一個裝修公司，效益不好關門了。而豬精也被騙，投資了一個所謂高科技塑膠廠，血本無歸。於是豬精努力轉了一下他的豬腦袋，用發現新大陸的語氣，對我們說教道：「高科技，都是騙人的！還是傳統行業好。」於是，他叫我姐姐去開店賣潮汕牛肉火鍋。至於我，作爲清華姚班畢業的高材生，應該學以致用，在淘寶開店賣牛肉丸。

雖然一直以來，我一直把豬精的話當豬叫，但是這段發言還是讓我下巴都驚掉了。但是更讓我震驚的是，我姐還真的聽了！在37虛歲的年紀，孑然一身，在潮汕牛肉火鍋店哼哧哼哧證明自己。每天給教授雞當狗，陪她逛街買包。還到處罵我不肯賣牛肉丸和陪教授雞逛街，實在是太不孝順了。

半年前，豬精房地產項目爛尾了，資金鏈斷裂。他聽說我自己賺錢了，於是跑到深圳我公司拜訪我，想跟我借錢。他對我說，他還是喜歡兒子多一點，但是因爲我賺到錢，也喜歡我。他說了一堆好聽話，然後送了我一塊2萬的手錶，再把我的100多萬的邁巴赫開走了。我姐姐聽說後，更痛苦了，她不明白，爲什麼她聽話賣牛肉火鍋，給雞當狗，做畜生的畜生，爲什麼爸爸還是愛我多過愛她。\footnote{我姐姐不知道李松堅借了我的邁巴赫拒絕歸還的事情，希望她知道我被李松堅爆金幣之後，能不那麼嫉妒我，開心一點。}

\ 

每次想到我從小相依爲命的姐姐一輩子給NPD豬精父親當血包的人生，我都忍不住流淚。我有時候會想，豬精如果不是這種10年以上，死緩以下的道德觀，也許我姐姐就不會這麼痛苦吧。如果父親不出軌，不犯故意傳播甲類傳染病罪、故意殺人未遂罪、濫用職權罪、組織黑社會性質組織罪等罪，我和姐姐現在應該正自豪地開着父親80年代買的豐田皇冠皇家沙龍，發抖音炫耀。而如果父親能突破犯罪上限，和凌菲菲一起把我姐姐肺結核搞死，一了百了，她也就不用忍受後來的痛苦。或者他能夠突破父女的倫理禁忌，也許他就會像愛狐狸精和雞一樣愛我姐姐，而我姐姐也不會一輩子缺愛了。

我從小就沒有爹，沒有感受過父愛，也因此也自由自在習慣了，實在不能理解我的姐姐爲什麼要給雞當狗，做畜生的畜生。或許是那個夏夜的父愛是那麼刻骨銘心，值得讓她終其一生苦苦追尋。

——亦或許，在那個夏夜，她腦子燒壞了吧。

\ \\
李新野\\
2025年7月19日

\ 

--------------------------------------------------------------------------

\

\textit{李新野網誌原文連結：https://sinyalee.com/blog/?p=1018}
